% SPDX-License-Identifier: GPL-3.0-or-later
\documentclass[11pt,a4paper]{article}

\usepackage[T1]{fontenc}
\usepackage[utf8]{inputenc}
\usepackage{lmodern}
\usepackage{microtype}
\usepackage[margin=25mm]{geometry}
\usepackage{amsmath,amssymb}
\usepackage{booktabs}
\usepackage{array}
\usepackage{graphicx}
\usepackage{caption}
\usepackage[numbers,sort&compress]{natbib}
\usepackage{xcolor}
\usepackage[colorlinks=true,linkcolor=blue!50!black,citecolor=blue!50!black,urlcolor=blue!50!black]{hyperref}

\newcommand{\Cp}{C_p}
\newcommand{\CNa}{C_{N\alpha}}
\newcommand{\CXzero}{C_{X0}}
\newcommand{\pinf}{p_\infty}
\newcommand{\fine}{\mathrm{fine}}
\newcommand{\tick}{yes}
\newcommand{\cross}{\textbf{no}}

\title{Verification and validation of the IgnisYeet aerodynamic pipeline\\
following the NASA WIND/NPARC procedure}
\author{IgnisYeet project}
\date{Draft, 17 October 2026}

\begin{document}
\maketitle

\begin{abstract}
IgnisYeet is a rocket simulator whose aerodynamic tool drives the open-source flow solver SU2.
We apply the verification and validation procedure of the NASA WIND/NPARC tutorial (AIAA G-077-1998)
to it in three layers: the solver alone against exact solutions (layer~0), the same cone through the
product's whole chain from CAD geometry to force coefficients (layer~1), and the tutorial's six steps on
a sample rocket and on the Basic Finner, including comparison with free-flight measurements (layer~2).
The aim is to attach a numerical uncertainty band to each coefficient and to set that band beside the
experimental error.

The results that are complete at this draft stage concern the solver and the method.
On four exact solutions (oblique shock, Prandtl--Meyer expansion, Taylor--Maccoll cone, Blasius plate),
SU2 run with the numerical settings the product writes reproduces the exact value to within 0.2\,\% on the
surface pressure and the shock angle in the primary setups, but wall quantities derived from the Mach number retain
errors of 0.8 to 3.6\,\% that decrease slowly or not at all.
Because the exact answer is known, we can test whether the grid convergence index (GCI) bounds the true
discretisation error: it does for 5 of 7 inviscid results and for 2 of 5 laminar plate results, and the two
plate cases that hold do so only where the error changes sign.
Neither the observed order nor the stability of that order across windows detects both failures, so a GCI
computed without an exact answer can understate the true error even when its order looks plausible.
Grids generated by the product cannot be refined similarly, because the mesher fixes the largest cell size;
a four-level Basic Finner study at Mach 2.375 does not establish a GCI, because the observed orders of its two
triplets disagree.
The NPARC archive's reference table for a 10$^\circ$ cone is internally inconsistent and disagrees with independent
tables in every entry except the surface Mach number.
% source: results/layer0.csv, results/layer0_gci.csv, results/gci_existing.csv, results/cone10_reference_check.csv (every number in this abstract)
% PENDING Task 18: abstract sentence on the layer 0 against layer 1 difference (pipeline error, including 3-D unstructured meshing)
% PENDING Task 14: abstract sentence on the sample-rocket Euler grid convergence and the refined d'Alembert check
% PENDING Task 15: abstract sentence on the Basic Finner RANS grid convergence at Mach 1.832 and the validation overlap judgement
% PENDING Task 16: abstract sentence on the iterative error relative to the grid uncertainty
% PENDING Task 20: abstract sentence on the Euler against RANS model difference relative to the numerical uncertainty
\end{abstract}

%=====================================================================
\section{Introduction}
\label{sec:intro}

Aerodynamic coefficients from a flow solver are numbers without a stated reliability unless they carry
an uncertainty.
The verification and validation procedure published by the NASA Glenn WIND/NPARC project
\citep{nparc_tutorial}, which follows the AIAA guide \citep{aiaa1998}, divides the task into six steps:
examination of the code, iterative convergence, consistency, spatial (grid) convergence, temporal convergence,
and comparison with a more accurate solution or with measurements.
Its central quantity is the grid convergence index of \citet{roache1998} in the form given by
\citet{celik2008}.
A comparison without an uncertainty band cannot say which side is wrong.

IgnisYeet is a rocket simulator.
Its aerodynamic tool takes a CAD description of the rocket, builds an unstructured volume mesh, runs SU2
\citep{economon2016} and returns force and moment coefficients to the flight simulator.
This paper applies the tutorial to that tool.
The deliverable is a numerical uncertainty band on the coefficients the tool returns, together with
an honest statement of where such a band can and cannot be computed.

\paragraph{Terminology.}
Three terms are kept apart throughout.
\emph{Verification} asks whether the equations are solved correctly and compares a result with an exact or
otherwise highly accurate solution.
\emph{Cross-comparison} compares one method or program with another; it shows that they differ or agree,
not which is closer to reality.
\emph{Validation} asks whether the model resembles the real system and compares a result with measured data.
Only the comparison with free-flight measurements in Section~\ref{sec:validation} is validation.
Layers~0 and~1 are verification.

\paragraph{Status of this draft.}
Several results are still being computed.
Passages that wait on them are marked in the source; their section headings appear in the typeset text without
content.
Each number rests on a committed result file or, where a comment in the source says so, on another committed file.

%=====================================================================
\section{Method}
\label{sec:method}

\subsection{Three layers}
\label{sec:layers}

% source: results/layer0.csv (header: SU2 version)
The study has three layers (Table~\ref{tab:layers}).
Layers~0 and~1 solve the same exact problem, a cone of 10$^\circ$ half angle in supersonic flow, by different
routes.
Layer~0 runs SU2~8.3.0 directly.
Layer~1 runs the product's whole chain.
Both are compared with the same Taylor--Maccoll solution, and the difference of their errors is the part due
to the chain around the solver.

The product accepts only a rocket geometry; a flat plate or a wedge cannot be built through it.
Layer~0 therefore does not exercise the product's pipeline, and no statement that the pipeline was verified
rests on layer~0.
Layer~0 uses two-dimensional and axisymmetric structured grids; layers~1 and~2 use three-dimensional unstructured
half models.
The difference between layers~0 and~1 thus includes the effect of three-dimensional unstructured meshing and
is reported as \emph{pipeline error, including three-dimensional unstructured meshing}, not as meshing error
alone.

\begin{table}[t]
\centering
\caption{The three layers of the study.}
\label{tab:layers}
\small
\begin{tabular}{@{}cp{4.4cm}p{4.2cm}p{3.6cm}@{}}
\toprule
Layer & What is run & What it tests & Kind of comparison \\
\midrule
0 & SU2 8.3.0 with the product's numerical settings, structured 2-D and axisymmetric grids
  & Solver and settings against exact solutions & Verification \\
1 & The product's chain, CAD to mesh to SU2 to coefficients, on the layer~0 cone
  & Chain as a whole, against the same exact solution & Verification \\
2 & The product on a sample rocket and on the Basic Finner
  & Error at working conditions; step~6 against free-flight data
  & Verification (steps 1--5), cross-comparison (Euler against RANS), validation (step~6) \\
\bottomrule
\end{tabular}
\end{table}

\subsection{The six steps}
\label{sec:steps}

Table~\ref{tab:steps} lists where each of the tutorial's steps is treated.
Step~5, temporal convergence, does not apply to the steady CFD of the product.
In its place the flight simulator's own time integrator is examined by step-size refinement
(Section~\ref{sec:temporal}).
Step~6 asks for a comparison with several turbulence models.
The product offers only the inviscid Euler model and Reynolds-averaged Navier--Stokes with the Spalart--Allmaras
model \citep{spalart1992}, so this part of the step is answered only in part (Section~\ref{sec:model}).

\begin{table}[t]
\centering
\caption{Tutorial steps and where they are treated.}
\label{tab:steps}
\small
\begin{tabular}{@{}cll@{}}
\toprule
Step & Content & Section \\
\midrule
1 & Examination of the code & \ref{sec:step1} \\
2 & Iterative convergence & \ref{sec:iter} \\
3 & Consistency & \ref{sec:consistency} \\
4 & Spatial (grid) convergence & \ref{sec:grid} \\
5 & Temporal convergence & \ref{sec:temporal} (integrator only) \\
6 & Model uncertainty and comparison with measurements & \ref{sec:model}, \ref{sec:validation} \\
\bottomrule
\end{tabular}
\end{table}

\subsection{Grid convergence index}
\label{sec:gci}

For three solutions $f_1,f_2,f_3$ on grids of representative size $h_1<h_2<h_3$, finest first, with
refinement ratios $r_{21}=h_2/h_1$ and $r_{32}=h_3/h_2$, let $\varepsilon_{21}=f_2-f_1$ and
$\varepsilon_{32}=f_3-f_2$.
For a constant ratio $r$ the tutorial gives the observed order
$p=\ln\!\big[(f_3-f_2)/(f_2-f_1)\big]/\ln r$, the extrapolated value $f_{h0}=f_1+(f_1-f_2)/(r^p-1)$ and
$\mathrm{GCI}=F_s\,|\varepsilon|/(r^p-1)$.
Where the ratio varies we use the procedure of \citet{celik2008}: with $s=\operatorname{sign}(\varepsilon_{32}/\varepsilon_{21})$,
\begin{equation}
p=\frac{1}{\ln r_{21}}\Big|\ln\Big|\frac{\varepsilon_{32}}{\varepsilon_{21}}\Big|+q(p)\Big|,\qquad
q(p)=\ln\frac{r_{21}^{p}-s}{r_{32}^{p}-s},
\end{equation}
solved by fixed-point iteration, then
$f_\mathrm{ext}=(r_{21}^{p}f_1-f_2)/(r_{21}^{p}-1)$ and
$\mathrm{GCI}_\fine=F_s\,e_a/(r_{21}^{p}-1)$ with $e_a=|(f_1-f_2)/f_1|$.
The factor of safety is $F_s=1.25$ for three grids and $3$ for two.
The asymptotic-range check is $r_{21}^{p}\,\mathrm{GCI}_{21}/\mathrm{GCI}_{32}$, which is near one in the
asymptotic range.
The grid size is $h=N^{-1/3}$ from the node count $N$ for three-dimensional grids and $h=N_\mathrm{cells}^{-1/2}$ for the
two-dimensional layer~0 grids, which are refined exactly by a factor of two in each direction.

% source: results/gci_existing.csv, results/layer0_gci.csv (headers: the gate, the order flags and the oscillatory rule)
Three rules decide when a number is reported.
\begin{itemize}
\item A grid level enters a study if the solver reported convergence or, where it gave no verdict, the residual fell by at
  least three orders of magnitude; a level the solver reported as not converged, or one with neither piece of
  evidence, is refused, and the refusal states the reason.
\item A computed order above 4 (twice the formal second order) or below 0.5 (a quarter of it) is flagged: the
  triplet is then outside the asymptotic range and its GCI is not reliable.
\item A triplet whose differences change sign (oscillatory convergence, $\varepsilon_{32}/\varepsilon_{21}<0$) is
  labelled and receives \emph{no} uncertainty band.
  \citet{celik2008} do compute a value for such a case; this study does not, and reports none.
\end{itemize}
Where the refinement is not geometrically similar, so that no single definition of $h$ is justified, the GCI
is computed under every defensible definition and both results are shown (Section~\ref{sec:grid}).

\paragraph{Acceptance of the procedure.}
% source: tests/test_gci.py, tests/data/prD.out (acceptance tests of the GCI procedure; not a results CSV)
The implementation reproduces the worked example of the tutorial \citep{nparc_spatial}
(three grids, $r=2$): it matches the printed output of the tutorial's own program, namely the order, the
extrapolated value, both GCI values and the asymptotic ratio, to a relative $10^{-4}$.
That program is single precision, so its fifth digit differs from a double-precision evaluation.
The double-precision values that the tests also pin are the study's own numbers, a regression check and not an
independent published reference.
Because that example has $r_{21}=r_{32}$, it does not exercise the variable-ratio path.
The procedure is therefore also tested against two monotone columns of Table~1 of \citet{celik2008}, which
have unequal ratios, to the precision that table prints.
One cell of that table, the extrapolated relative error in the second column, is printed as a value that
differs from the value recomputed from the same table's own printed entries; the cell is not used in the test.
The third column of that table is oscillatory; the implementation refuses it by the rule above rather than
reproduce the table's value.

\subsection{Exact solutions}
\label{sec:exact}

The exact solutions are computed, not transcribed: the weak-solution oblique shock relation, the
Taylor--Maccoll cone equation integrated from the shock to the surface with the shock angle found by shooting,
the Prandtl--Meyer function, and the Blasius similarity solution \citep{blasius1908,schlichting2017},
all for a perfect gas with $\gamma=1.4$.
The oblique-shock and Prandtl--Meyer results agree with the tabulated values of the NPARC archive pages
\citep{nparc_wedge,nparc_pm15} within the precision of those tables.
The cone solution is checked against independent tables in Section~\ref{sec:cone10}.
The Taylor--Maccoll solver in the product's documentation brackets the shock angle by a fixed upper limit that, for
large cone angles near detachment, lies on the wrong side of the maximum of the cone-angle curve; the 10$^\circ$
cones used here are far from that range, and the solver used in this study treats that case separately.
% source: no results CSV (comparison of the two solvers; not used for any number in this paper)
% source: results/layer0.csv (header: exact references and conditions)
The cases follow the NPARC archive \citep{nparc_archive}: a 15$^\circ$ wedge and a 15$^\circ$ expansion at Mach 2.5, a 10$^\circ$ half-angle cone at
Mach 2.35, and a laminar flat plate.
The plate is run at Mach 0.2 rather than the 0.1 of the archive case because the product's convective scheme has
no low-Mach preconditioning, so that a difference from Blasius is not confounded with low-Mach discretisation error.
The Blasius solution is incompressible, so a compressibility correction is estimated and reported alongside.

\subsection{Layer 0 configuration and grids}
\label{sec:l0setup}

% source: results/layer0.csv (header), results/layer0_gci.csv (header), results/layer0_plate.csv (header)
The numerical settings of the layer~0 Euler runs (convective discretisation, reconstruction, limiter, gradient
method, CFL control, linear solver and stopping criteria) are copied from a configuration that the product itself
wrote for an Euler case.
A test asserts that they are identical to the product's.
The only permitted differences are the boundary markers and, for the cone, the axisymmetric switch; the dimension
follows from the mesh.
The gas model is the product's.

The product has no laminar mode.
The laminar plate is run with the Navier--Stokes equations and without a turbulence model, using the flow
numerics and gas model of a product-written Reynolds-averaged configuration.
It therefore checks the product's RANS discretisation with the turbulence model removed;
it is not a configuration the product itself writes.
The product starts a RANS case with a first-order stage and restarts at second order; the layer~0 plate is
started cold at second order.
That the steady solution does not depend on this path was not tested.

The grids are body-fitted and structured, with three levels of $65\times49$, $129\times97$ and $257\times193$
nodes ($85\times49$, $169\times97$, $337\times193$ for the plate), exactly nested so that every coarse node is a
fine node and $r=2$ holds exactly.
In the supersonic cases the inlet, outlet and top are far-field boundaries, and every wave leaves through the
outlet: along the top boundary $|p/\pinf-1|\le7.3\times10^{-10}$ in each of the twelve runs.
The runs are driven to a density residual of $10^{-14}$ or to an iteration budget.
The convergence measure reported below is the peak minus the final value of $\log_{10}$ of the density residual.

\subsection{Layers 1 and 2}
\label{sec:l12setup}

Layer~1 solves the layer~0 cone, Mach 2.35, zero incidence, Euler, as a 10$^\circ$ cone-cylinder through the product
on a three-dimensional unstructured half model at three levels of resolution.
% PENDING Task 18: the realised layer 1 grids (cell sizes, node counts), the surface-pressure extraction window, and how the cone Cp is compared with Taylor-Maccoll

% source: results/gci_existing.csv (Basic Finner Euler conditions and levels), results/consistency.csv (sample-rocket Mach numbers)
Layer~2 treats a sample rocket (Euler, Mach 0.3, 0.8 and 2.0, incidence 0$^\circ$ and 4$^\circ$) and the
Basic Finner, a cone-cylinder body with four rectangular fins.
The Basic Finner is run with Euler at Mach 2.375 (incidence 0$^\circ$, 2$^\circ$, 4$^\circ$; four existing
levels) and with RANS (Spalart--Allmaras) at Mach 1.832 (incidence 0$^\circ$ and 2$^\circ$).
Mach 1.832 is chosen for RANS because it is the one point at which the product's RANS converges and at which
free-flight data exist; at Mach 2.375 the product's RANS stalls in a limit cycle and at Mach 3.147 it diverges.
% source: the product's validation documentation (not a results CSV)
The normal-force slope $\CNa$ is $C_N(2^\circ)/2^\circ$ for the Basic Finner, following the existing convention.
% PENDING Task 14: the realised sample-rocket Euler grid levels, machine and residual drops, and how CNalpha is formed there
% PENDING Task 15: the realised Basic Finner RANS levels and their residual drops

\paragraph{Grids built by the product cannot be similar.}
% source: results/gci_existing.csv (header and n1..n3 columns of basic_finner_mesh_check)
Refining the product's two cell-size settings together refines only the near-wall region.
The mesher grades the volume mesh from the wall cell size up to a largest cell size that is one tenth of the
far-field radius; that largest size is not a user setting and does not scale with the wall size.
Scaling it would require changing the far-field radius, which changes the problem solved.
For the four existing Basic Finner Euler levels, where wall and fin sizes scale by 1.4 per level, the cell count
grows by 1.70, 1.77 and 1.83 per level against $1.4^3=2.74$ for a similar refinement.
Neither $h\propto N^{-1/3}$ nor $h\propto$ wall size is then the effective grid size.
We compute every GCI of such a family under both definitions and treat the spread between them as part of the
uncertainty; the two definitions bracket the question and the spread is not a bound on the true GCI.

%=====================================================================
\section{Verification of the solver (layer 0)}
\label{sec:verification}

\subsection{A reference table that disagrees with itself}
\label{sec:cone10}

The NPARC archive offers a 10$^\circ$ cone at Mach 2.35 \citep{nparc_cone10}.
Its tabulated surface Mach number agrees with the Taylor--Maccoll value computed here to about $3\times10^{-3}\,\%$,
but its shock angle and its pressure, temperature and density ratios do not (Table~\ref{tab:cone10}).
The archive's shock angle is also inconsistent with its own surface Mach number and its own printed ratios.
We therefore do not use the archive's shock angle or ratios as targets; for this cone only the surface Mach number
and the surface pressure coefficient are compared.

\begin{table}[t]
\centering
\caption{NPARC cone10 table against the Taylor--Maccoll solution computed here (10$^\circ$ half angle, Mach 2.35, $\gamma=1.4$).}
\label{tab:cone10}
% source: results/cone10_reference_check.csv
\small
\begin{tabular}{@{}lrrr@{}}
\toprule
Quantity & NPARC & Taylor--Maccoll & Difference (\%) \\
\midrule
Shock angle $\beta$ (deg) & 27.1843 & 26.7367 & $-1.65$ \\
Surface Mach number & 2.1469 & 2.1468 & $-0.003$ \\
$p/\pinf$ & 1.4234 & 1.3739 & $-3.48$ \\
$T/T_\infty$ & 1.1063 & 1.0951 & $-1.01$ \\
$\rho/\rho_\infty$ & 1.2867 & 1.2546 & $-2.49$ \\
\bottomrule
\end{tabular}
\end{table}

To decide which shock angle is right, three further sources were collected (Table~\ref{tab:beta}).
Tabulated cone solutions \citep{sims1964}, interpolated to Mach 2.35, give $26.7366^\circ\pm0.02^\circ$.
The charts of \citet{ames1953}, drawn for $\gamma=1.405$ and read from the scan, give $26.80^\circ\pm0.15^\circ$,
uncorrected for a reading bias of about $+0.07^\circ$ that appears when the same method is applied to the tables.
A separate Taylor--Maccoll implementation inside the product gives $26.7367^\circ$.
None supports $27.18^\circ$.
The tables and the charts are independent of our code and of the archive page, in data and in implementation, but
they rest on the same theory: they solve the same Taylor--Maccoll equations for the same perfect gas.
They therefore test our implementation and the archive's number, not the theory.
The product's own implementation is only partly independent of ours: it can reveal a porting or numerical
error but not a shared conceptual one.

\begin{table}[t]
\centering
\caption{Shock angle of the 10$^\circ$ cone at Mach 2.35 from five sources. Uncertainty is the half-width of interpolation and reading error, not a statistical interval.}
\label{tab:beta}
% source: results/cone_shock_angle_sources.csv
\small
\begin{tabular}{@{}lrl@{}}
\toprule
Source & $\beta$ (deg) & Note \\
\midrule
This study's solver & 26.7367 & computed \\
NPARC cone10 page \citep{nparc_cone10} & 27.1843 & table text \\
Sims, SP-3004 Table 11 \citep{sims1964} & $26.7366\pm0.02$ & interpolated from tabulated Mach numbers \\
Ames, NACA 1135 Chart 5 \citep{ames1953} & $26.80\pm0.15$ & read from chart, $\gamma=1.405$ \\
Product's own implementation & $26.7367\pm0.001$ & same equations, separate code \\
\bottomrule
\end{tabular}
\end{table}

\subsection{Wedge, cone and expansion}
\label{sec:l0results}

% source: results/layer0.csv, results/layer0_gci.csv
All twelve layer~0 Euler runs reached the residual target (12.7 to 13.9 orders of density-residual drop, 132 to
230 iterations).
Table~\ref{tab:l0} gives the finest-grid error against the exact solution, the observed order, the GCI and
whether the GCI bounded the true finest-grid error.

\begin{table}[t]
\centering
\caption{Layer~0 Euler results on the finest grid ($257\times193$ nodes) against the exact solution, with the grid convergence index of the three-grid sequence.
``Bounds'' says whether $\mathrm{GCI}_\fine\ge|\text{error}|$.
The order of a quantity marked $^\dagger$ is below 0.5 and its GCI is flagged as unreliable. Wall quantities are means over a window of wall nodes.}
\label{tab:l0}
% source: results/layer0_gci.csv
\footnotesize
\setlength{\tabcolsep}{4pt}
\begin{tabular}{@{}llrrrrrc@{}}
\toprule
Case & Quantity & Exact & Finest & Error (\%) & $p$ & $\mathrm{GCI}_\fine$ (\%) & Bounds \\
\midrule
Wedge, M 2.5 & wall $p/\pinf$ & 2.4675 & 2.4676 & 0.0059 & 0.82 & 0.011 & \tick \\
 & shock angle (deg) & 36.945 & 36.959 & 0.037 & 0.68 & 0.099 & \tick \\
Cone, M 2.35 & wall $\Cp$ & 0.09673 & 0.09655 & $-0.188$ & 3.39 & 0.059 & \cross \\
 & wall Mach & 2.1468 & 2.1297 & $-0.80$ & 0.39$^\dagger$ & 2.48 & \tick \\
Cone, symmetry-axis variant & wall $\Cp$ & 0.09673 & 0.09601 & $-0.75$ & 1.49 & 0.71 & \cross \\
 & wall Mach & 2.1468 & 2.0828 & $-2.98$ & 0.77 & 4.50 & \tick \\
Expansion, M 2.5 & wall $p/\pinf$ & 0.32743 & 0.32746 & 0.0097 & 0.46$^\dagger$ & 0.077 & \tick \\
 & wall Mach & 3.2368 & 3.1218 & $-3.55$ & none & none & -- \\
\bottomrule
\end{tabular}
\end{table}

Two features of the table need to be read with care.
First, the observed orders range from 0.39 to 3.39, against a formal second-order scheme.
They are not a measurement of the scheme's order.
The wall windows contain the tip or corner transient, and the shock angle is read from a captured, smeared
shock; the notes of the result file describe the low orders as consistent with these causes rather than with a
smooth-field discretisation error.
For the wedge wall pressure the order is stable across four moving windows (0.84 to 1.04); for the expansion wall
pressure it is not (0.65 to 2.54).
Second, the wall Mach number on the expansion does not converge: the three grids give 3.1214, 3.1216 and 3.1218
against an exact 3.2368, an error of $-3.57$, $-3.56$ and $-3.55\,\%$, essentially unchanged by refinement, so no order or GCI
exists although the run itself converged (13.1 orders on the finest grid).
A GCI is a statement about discretisation error that shrinks with the grid, and this error does not.

\paragraph{Wall entropy.}
% source: results/layer0.csv (diag_entropy_* rows and header)
Along the wall the entropy exceeds the exact value on all three cases, while ten rows off the wall it is at the
exact value and upstream of the corner or tip it is exact.
On the expansion the wall entropy ratio $s/s_\infty$ is 1.0785, 1.0784 and 1.0783 on the three grids, a constant
offset of about 7.8\,\%.
On the wedge the wall reads 1.058 on all three grids, ten rows off the wall 1.030 on the two finer grids, and the exact post-shock
value is 1.030: a grid-independent excess of 2.7\,\% at the wall.
On the cone the excess falls with the grid (3.76, 1.98 and 1.06\,\% in the surface window).
This is consistent with entropy generated at the corner, leading edge or tip and convected along the wall
streamline; that is an interpretation, not a proof.
The surface pressure on the expansion has an error of 0.0097\,\% while the wall Mach number has $-3.55\,\%$, so
quantities obtained from the wall pressure show no comparable error, whereas quantities taken from the wall Mach
number, density or temperature would be expected to carry the excess.

\paragraph{The tip junction.}
% source: results/layer0.csv (cone and cone_sym_axis rows), results/layer0_gci.csv (header)
The cone is run twice.
In the first run the symmetry axis and the cone surface form one slip-wall boundary; in the second the axis is a
symmetry boundary and the cone surface a slip wall.
In the second run the node at the tip belongs to both boundaries and is a stagnation node on every grid (Mach 0,
$p=2.76\,\pinf$), whereas in the first run it is regular (Mach 2.23, $p=1.17\,\pinf$).
The finest-grid error of the wall pressure coefficient is $-0.19\,\%$ in the first run and $-0.75\,\%$ in the
second, and of the wall Mach number $-0.80$ and $-2.98\,\%$.
The comparison of the two runs shows that most of the second run's error is the boundary-condition junction.
The first run is taken as the layer~0 cone result.
On the expansion, splitting the lower boundary into two slip-wall boundaries at the corner made the corner node a
stagnation point and stalled the density residual near $-1.5$ over 20\,000 iterations.
These are sensitivities of a two-dimensional structured setup.
The product's three-dimensional meshes also have boundary junctions (wall to base), but whether they are affected
in the same way was not tested.

\subsection{Does the GCI bound the true error?}
\label{sec:gcibound}

% source: results/layer0_gci.csv (true_error_pct, gci_bounds_error; header)
The purpose of a GCI is to bound the discretisation error.
Layer~0 is the only place in this study where the exact answer is known, so it can test whether it does.
Of the seven inviscid rows with a computed GCI, the bound held for five and failed for two:
the cone wall pressure coefficient (GCI 0.059\,\% against a finest-grid error of 0.188\,\%, so the true error is
3.2 times the GCI) and the same quantity in the symmetry-axis variant (0.71\,\% against 0.75\,\%).
One further row, the expansion wall Mach number, has no GCI.

``Held'' does not mean tight.
The ratio of GCI to true error is about 8 for the expansion wall pressure, 3.1 for the cone wall Mach number and
1.5 for the symmetry-axis cone wall Mach number (ratios of the entries of Table~\ref{tab:l0}).
A bound that is eight times the error is a bound, not an estimate.
Rows flagged as outside the asymptotic range (orders below 0.5 or above 4) are included in the count of five.

The practical question is whether a failure can be detected when no exact answer exists, which is the situation
in layer~2 (Table~\ref{tab:detect}).
The order flag (outside 0.5 to 4) misses both failures, because the cone wall pressure coefficient has $p=3.39$.
Stability of the order across four moving windows (within 0.5) catches the cone failure, whose orders range from
2.10 to 4.12, but misses the symmetry-axis failure, whose orders are stable at 1.36 to 1.58;
and it flags the expansion wall pressure as unstable (0.65 to 2.54) although the GCI bounded that error.
Neither available check is therefore sufficient, and a GCI computed without an exact answer can understate the
true error even when its order looks plausible and its windows look stable.

\begin{table}[t]
\centering
\caption{Can a failure of the GCI be detected without the exact answer? Rows are the two failures and one held case for contrast. Orders over four windows $[0.3,1]$, $[0.4,1]$, $[0.5,1]$, $[0.6,1]$ of the wall; stable means within 0.5 of each other.}
\label{tab:detect}
% source: results/layer0_gci.csv (order, note, gci_bounds_error)
\small
\begin{tabular}{@{}lccccc@{}}
\toprule
Quantity & $p$ & Order flag? & Window orders & Stable? & GCI bounds? \\
\midrule
Cone wall $\Cp$ & 3.39 & no & 2.10--4.12 & no & \cross \\
Symmetry-axis cone wall $\Cp$ & 1.49 & no & 1.36--1.58 & yes & \cross \\
Expansion wall $p/\pinf$ & 0.46 & yes & 0.65--2.54 & no & \tick \\
\bottomrule
\end{tabular}
\end{table}

\subsection{Laminar flat plate against Blasius}
\label{sec:plate}

% source: results/layer0_plate.csv, results/layer0_gci.csv (plate rows)
The laminar plate at Mach 0.2 ($Re=10^5$ per metre) is solved on three nested grids of $84\times48$, $168\times96$
and $336\times192$ cells (see Section~\ref{sec:l0setup} for the configuration).
All three runs reached the residual target, with a density-residual drop of 9.4 to 10.4 orders.
Table~\ref{tab:plate} gives the differences from the Blasius solution of the skin-friction coefficient at $x=0.5$
and of the displacement and momentum thicknesses of the profile there.

\begin{table}[t]
\centering
\caption{Laminar plate at Mach 0.2, differences from Blasius (\%) on the coarse, medium and fine grid, the difference of the Richardson extrapolation from Blasius, observed order and GCI. ``Scaled'' quantities use the profile's own edge velocity; those rows were added after the unscaled momentum thickness failed to be bounded, so they are post hoc, and the plate rows are not independent (each thickness appears scaled and unscaled).}
\label{tab:plate}
% source: results/layer0_gci.csv (plate rows), results/layer0_plate.csv
\footnotesize
\setlength{\tabcolsep}{4pt}
\begin{tabular}{@{}lrrrrrrr@{}}
\toprule
Quantity & Coarse & Medium & Fine & Extrapolated & $p$ & $\mathrm{GCI}_\fine$ (\%) & Bounds \\
\midrule
$C_f$ at $x=0.5$ & $-1.95$ & $-0.71$ & $+0.008$ & $+0.97$ & 0.80 & 1.20 & yes, sign crossing \\
$\delta^*$ & $+2.01$ & $+0.37$ & $-0.42$ & $-1.13$ & 1.07 & 0.89 & yes, sign crossing \\
$\theta$ & $+0.46$ & $-1.13$ & $-1.87$ & $-2.50$ & 1.11 & 0.80 & \cross \\
$\delta^*$, scaled & $+3.52$ & $+1.86$ & $+1.07$ & $+0.33$ & 1.05 & 0.92 & \cross \\
$\theta$, scaled & $+3.33$ & $+1.73$ & $+0.98$ & $+0.30$ & 1.08 & 0.84 & \cross \\
\bottomrule
\end{tabular}
\end{table}

The bound held for two of five rows and failed for three.
The two that held do so only because the sequence crosses Blasius between the medium and the fine grid:
the fine-grid $C_f$ error of $+0.008\,\%$ is a sign crossing ($-1.95$, $-0.71$, $+0.008\,\%$), not agreement, and a
finest-grid error smaller than the GCI there is not evidence that the GCI bounds the error of a converged solution.
What remains after extrapolation is $+0.97\,\%$ on $C_f$, which the GCI cannot account for.

The difference from Blasius mixes several things that a discretisation-error bound does not cover.
(i)~The outer flow at $\eta=8$ is $1.0036\,U$ on every grid, and for a boundary layer under an edge velocity
$u_e$ the skin friction scales as $(u_e/U)^{3/2}$; $C_f/(u_e/U)^{3/2}$ on the fine grid is $0.53\,\%$ below
Blasius (an estimate that takes $u_e$ at one height).
(ii)~An Eckert reference-temperature estimate of compressibility at Mach 0.2 changes $C_f$ by $-0.054\,\%$
(an estimate; the CFD value is not corrected).
(iii)~The domain matters.
In a first, shorter domain, with the plate ending at the outlet, $C_f$ on the medium grid was $+3.65\,\%$ from the
value in the domain used here; two other extents of the domain changed it by $+0.30\,\%$ and $-0.16\,\%$.
The domain used here was chosen after seeing the short-domain result, not before.
Two further enlargements were tried in exploratory runs not used here, and neither met the residual target on the
medium grid; the domain is not shown converged.
(iv)~The leading-edge node, which belongs to both the symmetry plane and the wall, is a stagnation node with a pressure spike whose
coefficient grows with refinement (0.022, 0.032, 0.047 on the three grids) and does not converge;
no single-boundary remedy is available because a no-slip wall and a symmetry plane cannot be one boundary.
(v)~For each quantity the observed order is the same within 0.12 at the stations $x=0.25$ to 0.75 (0.75 to 1.16 over
all quantities), below the formal second order, so it cannot separate the order of the scheme from a leading-edge error that the boundary
layer carries downstream.
Moreover $C_f\sqrt{Re_x}$ on the finest grid is 0.6591, 0.6641, 0.6663 and 0.6680 at $x=0.25$, 0.5, 0.75 and 1
against Blasius 0.664: it rises through the exact value with $x$ and does not approach it from one side.
What sets that slope was not separated.

The laminar plate therefore does not give a clean verification of the discretisation.
On the finest grid the differences from Blasius are $+0.008\,\%$ on $C_f$ (at a sign crossing) and 0.4 to 1.9\,\% on
the profile integrals; the remaining difference is attributable in part, but not fully quantified, to the effects
(i)--(v).

\subsection{Turbulent flat plate: grid convergence of existing data}
\label{sec:turbplate}

% source: results/gci_existing.csv (rans_flatplate rows)
Two existing three-grid Reynolds-averaged plate studies (Mach 0.2, the flat-plate case of the NASA Turbulence
Modeling Resource \citep{tmr}, and Mach 2; grids of 3\,264, 13\,056 and
52\,224 cells, $r=2$) were analysed for the skin-friction coefficient at $x=0.97$.
The observed orders are 1.36 (Mach 0.2) and 1.12 (Mach 2), the GCI of the finest grid is 0.038\,\% and 0.079\,\%,
and the asymptotic ratios are 0.9995 and 0.9993.
These are grid-convergence statements only: no exact solution is used, so nothing here tests whether the GCI
bounds the error.
Orders of 1.1 to 1.4 lie below the formal order of 2 and are not flagged by the order rules; a limiter is a
possible cause that was not tested.
At Mach 2 the medium grid stopped at the 20\,000-iteration cap with a residual drop of 5.76 orders, against 7.81 to
8.01 on the other grids; it passes the three-order gate and missed its own stopping target.
The configuration that produced these data was written by a script of the product's documentation, not by the
product.
It differs from the product's second-order configuration in its stop rule, its start-up stage and its gas constant,
among other things, and these differences are pinned by a test;
so what was verified is a configuration close to, but not identical to, the product's.
% source: tests/test_rans_flatplate_script.py (test of the differences; not a results CSV)

\subsection{Conservation of mass}
\label{sec:massflux}

% source: results/massflux.csv
SU2 reports the mass flux through each boundary marker.
On the planar cases (wedge, expansion) the inflow and outflow agree to the print floor of the history file
(relative imbalance at or below $10^{-9}$) on every grid.
That is what a conservative scheme gives at a converged state, so it tests convergence and flux bookkeeping, not the
discretisation error.
On the axisymmetric cone the flux is for the full revolution and carries the $2\pi r$ weight; the imbalance is
nonzero on every grid, $6.6\times10^{-4}$, $1.8\times10^{-4}$ and $4.3\times10^{-5}$ on the three grids
(successive ratios 3.75 and 4.07), and $3.4\times10^{-4}$, $1.0\times10^{-4}$ and $2.6\times10^{-5}$ in the
symmetry-axis variant (ratios 3.33 and 3.99).
Every value is below the tutorial's $10^{-3}$ for a consistent solution.
A ratio near four per refinement is consistent with a second-order discretisation error; with three grids it is an
observation and not an observed order.
That the $2\pi r$ weight is the cause is an interpretation consistent with these numbers, not a proof.
The symmetry-axis variant has 0.52 to 0.59 times the imbalance of the first, a difference that is not explained.

\subsection{Layer 1: the cone through the product}
\label{sec:layer1}

% PENDING Task 18: layer 1 cone-cylinder results: surface Cp against Taylor-Maccoll on three levels, shock angle, cone-only axial force, GCI per quantity (or the refusal with its residual drop), grid-similarity caveat
% PENDING Task 18: the layer 0 against layer 1 comparison table and the pipeline contribution, called pipeline error including 3-D unstructured meshing; the statement of whether the two layers are distinguishable at this resolution given each layer's GCI

%=====================================================================
\section{Application: the six steps on the sample rocket and the Basic Finner (layer 2)}
\label{sec:application}

\subsection{Step 1: examination of the code}
\label{sec:step1}

% source: no results CSV; the product's test documentation and the tests added by this study
No new calculation is made for this step.
The product's existing tests include order tests of its time integrators.
Its tests do not measure the order of mesh convergence, and no test of grid-convergence code existed before this
study; the CFD solver itself is not run by the tests that run in continuous integration.
This study adds tests for the GCI procedure (Section~\ref{sec:gci}), for the exact solutions against independent
tables (Section~\ref{sec:cone10}), for the identity of the layer~0 numerical settings with the product's, and for
the consistency checks below.

\subsection{Step 2: iterative convergence}
\label{sec:iter}

The tutorial asks for three things: that the residual falls by several orders, that the coefficients become flat,
and that what remains after the iteration is stopped is small against the grid uncertainty.
The last is what decides whether a grid study is meaningful; if the iterative error is of the order of the GCI,
the grid convergence argument fails.
% source: results/layer0.csv, results/layer0_plate.csv (residual_drop_orders)
In layer~0 every run reached its residual target, with 12.7 to 13.9 orders of drop for the Euler cases and 9.4 to 10.4
for the three plate grids (Sections~\ref{sec:l0results} and~\ref{sec:plate}).
% PENDING Task 16: the iterative-error to GCI ratio per layer 2 condition, with any condition above the thresholds stated plainly

\subsection{Step 3: consistency}
\label{sec:consistency}

% source: results/consistency.csv
\paragraph{$C_N(\alpha=0)=0$ and antisymmetry $C_N(-\alpha)=-C_N(\alpha)$.}
% source: results/consistency.csv
For a symmetric body both relations follow from the symmetry of the shape.
They were checked on five configurations of an existing Euler asymmetry study at Mach 0.3, 0.8 and 2.0.
On the four configurations without mirror symmetry of the mesh (12 rows), $|C_N(0)|$ is 0.013 to 0.091, against
tolerances of 0.007 to 0.012 (1\,\% of the mirror-combined normal force), and the mirror-combined even part
$[C_N(+4^\circ)+C_N(-4^\circ)]/2$ is 0.014 to 0.085 in magnitude; every one of these 24 rows fails.
On the fifth configuration, a mesh made mirror-symmetric by construction, the checks pass, but they pass by construction
and say nothing about the solver.
The production runs of the sample rocket use such a symmetric mesh; their $C_N(0)$ is zero to solver round-off
($|C_N|<2\times10^{-5}$) for the same reason.
The tests cannot separate a solver effect from a mesh-asymmetry effect on the asymmetric meshes.
(Values in the source are rounded to five decimals, so a printed zero means below $5\times10^{-6}$.)
% PENDING Task 14: antisymmetry from directly solved positive and negative incidence on the production mesh

\paragraph{d'Alembert's paradox.}
% source: results/consistency.csv (ca_dalembert row)
At Mach 0.3 and zero incidence the inviscid axial force coefficient is expected to be zero (d'Alembert's paradox).
On the production grid of the sample rocket $C_A=-0.031$.
This is a single grid, so no statement is made here on whether it is a discretisation error that shrinks with the
mesh or a defect of formulation or boundary conditions.
% PENDING Task 14: the refined-grid |CA| at Mach 0.3, alpha 0 on three similar levels, and which of the two readings (shrinks with the grid, or does not) it supports

\paragraph{Mass-flux balance.}
The balance was checked on the layer~0 cases (Section~\ref{sec:massflux}); it was not checked on the layer~2 far
field.

\subsection{Step 4: spatial convergence}
\label{sec:grid}

\subsubsection{Reynolds-averaged sample-rocket family from existing data}

% source: results/gci_existing.csv (rans_mesh_study rows)
An existing three-level RANS family of the sample rocket (wall size 8, 4.8 and 3\,mm; Mach 0.8 and 2; incidence
$0^\circ$ and $\pm4^\circ$, the $4^\circ$ values combined from $+4^\circ$ and $-4^\circ$ solutions) was passed through
the convergence gate.
All four families are refused (Table~\ref{tab:refused}).
The study's data file records the solver's verdict for each solve but no residual drop, so a residual drop is
not available for any of these levels.

\begin{table}[t]
\centering
\caption{RANS sample-rocket families refused by the gate.}
\label{tab:refused}
% source: results/gci_existing.csv (rans_mesh_study rows, note column)
\small
\begin{tabular}{@{}llp{7.3cm}@{}}
\toprule
Mach & Quantities & Reason \\
\midrule
0.8 & $C_N$, $C_A$, $x_{cp}$ at $4^\circ$ & The 8\,mm level's $-4^\circ$ solve was recorded as not converged \\
0.8, 2 & $C_N(0)$ & The $0^\circ$ solve was recorded as not converged at 8\,mm and 4.8\,mm \\
2 & $C_N$, $C_A$, $x_{cp}$ at $4^\circ$ & The 3\,mm level's $-4^\circ$ solve is a rerun with a fixed iteration count and no recorded verdict or residual drop \\
\bottomrule
\end{tabular}
\end{table}

The product's own convergence verdict ``unconverged'' includes a case that reached the iteration limit with forces
and moments nearly constant over the last window, in which case the window average is used and a warning issued;
such a case has flat coefficients but not necessarily the residual drop that the tutorial's step~2 asks for.
% source: the product's CFD documentation (not a results CSV)
The consequence is stated narrowly.
The grid-convergence index published for this family at Mach 0.8 rests on a coarse-level value that combines a $-4^\circ$
solve recorded as not converged; it therefore does not meet the tutorial's iterative-convergence precondition on
the residual-drop leg (the coefficient-flatness leg may be met), and the documentation does not say so.
% source: the product's CFD documentation (not a results CSV)
This is evidence that the precondition is not demonstrated; it is not a claim that the published number is wrong.

\subsubsection{Basic Finner, Euler, Mach 2.375}

% source: results/gci_existing.csv (basic_finner_mesh_check rows, fairing)
Four existing levels (345\,736 to 1\,891\,652 cells) give two triplets, L1--L3 (coarse) and L2--L4 (fine), each
under both grid-size definitions (Table~\ref{tab:finner}).
The grids are not geometrically similar (Section~\ref{sec:l12setup}).
The source records a convergence status for each level but no residual drop.

\begin{table}[t]
\centering
\caption{Basic Finner, Euler, Mach 2.375, base treatment ``fairing'': observed order $p$ and $\mathrm{GCI}_\fine$ (\%) of two triplets under both definitions of the grid size ($h\propto N^{-1/3}$, $h\propto$ wall size). A dash means no order was computed: the differences change sign, or the finer pair changes at least as much as the coarser, so there is no GCI. Values marked $^\dagger$ have an order outside 0.5 to 4 and are flagged as unreliable.}
\label{tab:finner}
% source: results/gci_existing.csv (basic_finner_mesh_check, fairing)
\footnotesize
\setlength{\tabcolsep}{3.5pt}
\begin{tabular}{@{}llcrrrr@{}}
\toprule
Quantity & Incidence & Triplet & $p$ ($N^{-1/3}$) & GCI & $p$ (wall) & GCI \\
\midrule
$C_A$ pressure part & 0$^\circ$ & L1--L3 & 1.28 & 22.7 & 0.48$^\dagger$ & 35.2 \\
 & & L2--L4 & 2.99 & 4.35 & 1.58 & 4.99 \\
 & 2$^\circ$ & L1--L3 & 1.31 & 22.0 & 0.499$^\dagger$ & 33.8 \\
 & & L2--L4 & 2.85 & 4.74 & 1.50 & 5.46 \\
 & 4$^\circ$ & L1--L3 & 1.55 & 17.1 & 0.63 & 24.7 \\
 & & L2--L4 & 2.69 & 4.95 & 1.41 & 5.72 \\
$C_N$ & 2$^\circ$ & L1--L3 & 8.45$^\dagger$ & 0.034 & 4.33$^\dagger$ & 0.040 \\
 & & L2--L4 & -- & -- & -- & -- \\
 & 4$^\circ$ & L1--L3 & 0.52 & 2.62 & 0.07$^\dagger$ & 11.1 \\
 & & L2--L4 & 2.24 & 0.33 & 1.15 & 0.39 \\
$x_{cp}$ & 2$^\circ$ & L1--L3 & -- & -- & -- & -- \\
 & & L2--L4 & -- & -- & -- & -- \\
 & 4$^\circ$ & L1--L3 & -- & -- & -- & -- \\
 & & L2--L4 & -- & -- & -- & -- \\
\bottomrule
\end{tabular}
\end{table}

The two triplets disagree.
For the pressure part of the axial force the coarse triplet gives orders of 1.3 and 0.5 (the two definitions) and the fine one
3.0 and 1.6; for the normal force at $4^\circ$ the orders are 0.52 and 0.07 against 2.24 and 1.15.
The normal force at $2^\circ$ gives a very small coarse-triplet GCI at an order above 4, and the fine triplet does not
converge.
The centre-of-pressure position is oscillatory or non-converging in every triplet and receives no band.
The normal-force slope is formed from $C_N$ at $2^\circ$, for which the coarse triplet gives a flagged order (8.45 and 4.33) and the fine triplet
does not converge.
The only band available for the slope is therefore a sensitivity and not a GCI: the coarse triplet gives $\pm0.17\,\%$ for the
plotted grid (0.167 to 0.173\,\% over both definitions), while the largest observed change from the plotted grid to any
other level, L2 to L4, is $+0.47\,\%$ and is larger than that.
With the flat-base treatment, $C_N$ and $x_{cp}$ at $2^\circ$ and $4^\circ$ are non-converging or oscillatory in every triplet.
This four-level study does not establish a reliable GCI for the Basic Finner.
% source: results/gci_existing.csv, derived L2-to-L4 change in cn at fairing alpha 2: 0.33008/0.32852 - 1

\subsubsection{Sample rocket, Euler, similar family}

% PENDING Task 14: the three similar Euler levels of the sample rocket at Mach 0.3, 0.8 and 2.0, incidence 0 and 4 deg: node counts, residual drops, GCI of CNalpha, CA and xcp under both h definitions, the order flags, and any oscillatory entry stated as without a band

\subsubsection{Basic Finner, RANS, Mach 1.832}

% PENDING Task 15: the three RANS levels at Mach 1.832 (incidence 0 and 2 deg), residual drop per level, GCI of CNalpha, CA and xcp, or the statement that a RANS GCI was not obtained with the residual drop

\subsection{Step 5: temporal convergence}
\label{sec:temporal}

% source: results/temporal.csv
The CFD is steady, so temporal convergence in the tutorial's sense does not apply to it; there is no time step in the
SU2 runs.
The product's time-dependent part is the six-degree-of-freedom flight integrator, and its convergence is examined
by step-size refinement.
The sample rocket is flown with a fourth-order Runge--Kutta integrator with steps of 8, 4, 2 and 1\,ms
($r=2$; parachute descent, 4\,m/s wind), giving two triplets, 1/2/4\,ms and 2/4/8\,ms.
The quantities are apogee, apogee time, landing distance, flight time, rail-exit time and speed, and altitude and speed at
fixed times.

No order was obtained for the sample rocket.
Every row is refused: the differences change sign (oscillatory), or two steps give identical output at its printed
resolution.
Over the four step sizes the apogee altitude varies by $0.018$\,m in $4440.5$\,m ($4.1\times10^{-6}$ relative), the landing
distance by 0.032\,m in 7735.5\,m and the flight time by 2.4\,ms in 685\,s.
That is a spread, not an error estimate: without a convergence order no band is claimed.

The scatter is attributed, on the evidence of controls, to the rail model.
The rocket is held at rest on the rail while the thrust is below its weight, which puts a kink in the acceleration at
about 2.3\,ms that the integrator does not locate; the kink falls in a different step relative to the step ends in each run.
Three controls modify the motor thrust curve (not the product's sample):
removing both the hold and the 0 to 0.05\,s thrust ramp gives orders of 4.01 for apogee and 3.96 for flight time on
the coarse triplet, as expected of a fourth-order scheme;
removing only the hold reduces the rail-exit speed difference between the 8 and 4\,ms runs from $4.5\times10^{-4}$ to
$4.1\times10^{-8}$\,m/s, which supports the hold as the main cause of the rail-exit scatter;
and a fully nested sweep rules out the thrust-curve knots that the step does not divide as the dominant cause.
A second, smaller contribution is not excluded: with only the hold removed, the coarse-triplet apogee is still
non-monotone.
The finest triplets of the controls stop converging at about $10^{-9}$ to $10^{-8}$ relative, a floor that
persists on the nested sweep and whose cause is unidentified.
The one order from a finest triplet (4.75 for the apogee time of the first control) has differences below that
floor and the same run's apogee oscillates, so it is not evidence of fourth-order convergence.
This is verification of the integrator by step refinement; it is not a comparison with measured data.

\subsection{Step 6: model uncertainty}
\label{sec:model}

The product offers Euler and Reynolds-averaged Navier--Stokes with the Spalart--Allmaras model, and nothing else.
The tutorial's request for a comparison of several turbulence models is therefore answered only in part, and the
limitation is a limit of the product's scope, not an omission.
The layer~0 comparison of Spalart--Allmaras with a second turbulence model on the turbulent plate was not performed.
Comparing Euler with RANS-SA on the same condition is a cross-comparison of two models; a difference between
them counts as a model difference only if it exceeds the sum of the two numerical uncertainty bands.
% PENDING Task 20: the Euler against RANS-SA differences (sample rocket Mach 0.8 and 2.0 at 4 deg; Basic Finner Mach 1.832 at 0 and 2 deg), each with both GCI bands and the judgement whether the difference is separable from numerical uncertainty

%=====================================================================
\section{Validation against free-flight data}
\label{sec:validation}

% source: figures/validation_bands.svg (the plotted measured and CFD values are read from the Basic Finner data files of the product's documentation, which are not copied into results/)
The Basic Finner coefficients are compared with the free-flight measurements of \citet{dupuis1997}: aeroballistic
range tests, with the coefficients obtained from a six-degree-of-freedom maximum-likelihood fit
over groups of shots at eight Mach numbers between 1.077 and 4.300.
Each tabulated coefficient carries a probable error in per cent.
The error bars follow the table.
A blank entry means the error is unknown, and no bar is drawn; an error is never estimated and filled in.
A tabulated 0 means an error below 0.5\,\%, not zero error, and is drawn as a hollow marker with caps at the
0.5\,\% bound.
An existing validation figure of the product's documentation draws a bar of length zero for a quantity whose error is
unknown; the figure used here does not.
% source: the product's documentation figure code (not a results CSV)

\begin{figure}[t]
\centering
\includegraphics[width=\textwidth]{fig/validation_bands.pdf}
\caption{Basic Finner $\CNa$ (left) and $\CXzero$ (right) against Mach number: free-flight measurements with
tabulated probable error, Euler CFD, and the RANS points at Mach 1.832 (two treatments of the base term).
The hatched box on the Euler point at Mach 2.375 is a grid sensitivity ($\pm0.17\,\%$, too small to see at this scale), not an
established GCI.
There is no band on $\CXzero$ and none at any other Mach number.}
\label{fig:validation}
\end{figure}

Every CFD point in Fig.~\ref{fig:validation} is compared with the measured value only; no difference between CFD
models is drawn or computed.
The Euler $\CNa$ and the Euler $\CXzero$ lie below the measured value at every Mach number plotted.
Whether the difference between Euler and measurement can be attributed to model form rather than numerical error
cannot be decided: there is no established numerical uncertainty band on the Euler coefficients.
The grids are not geometrically similar, the orders of the two triplets disagree (Section~\ref{sec:grid}),
and the only band that can be drawn, $\pm0.17\,\%$ on $\CNa$ at Mach 2.375, is not a GCI and is smaller than the
0.47\,\% change observed between two of the grids.
No band is drawn on $\CXzero$ because the pressure part of the axial force, whose grid study is available, is only one
of several terms of $\CXzero$.
% PENDING Task 15: the figure regenerated with the RANS band, the RANS CNalpha at Mach 1.832 with its GCI band (or the statement that none was obtained), and the comparison with the measured value and its probable error
% PENDING Task 15: the judgement of whether the CFD band and the measured error bar overlap and, if they do not, the statement that the difference cannot be explained by numerical error
% PENDING Task 20: the validation statement on model uncertainty at Mach 1.832 (Euler against RANS-SA against measurement)

%=====================================================================
\section{Discussion}
\label{sec:discussion}

\paragraph{What the GCI can be trusted for.}
% source: results/layer0_gci.csv (counts, true_error_pct, gci_bounds_error; the ratio of 8 is derived from gci_fine_pct and true_error_pct of the expansion wall pressure)
The central result of layer~0 is a measurement of the GCI itself.
Where the exact answer is known it bounded the true error in 5 of 7 inviscid cases and 2 of 5 laminar-plate cases.
The two inviscid failures were not both flagged by the order rule or by the window-stability test (Table~\ref{tab:detect}).
Holding is not tightness: the bounds that held were up to eight times the error.
A GCI computed in layer~2, without an exact answer, is therefore a conditional estimate that can understate the true
error, and should be read with the flags of Section~\ref{sec:gci} and not as a certificate.

\paragraph{Wall quantities derived from the Mach number.}
% source: results/layer0.csv (expansion wall_mach, wall_p_ratio and diag_entropy_wall_window rows)
A wall entropy excess that does not change with the grid accompanies a wall Mach error of 3.55\,\% on the expansion that does not
decrease with refinement and for which no GCI exists.
Surface pressure, from which forces are integrated, has an error of 0.01\,\% there.
The interpretation (entropy generated at a corner or tip and convected along the wall) is consistent with the data,
not proven.

\paragraph{Boundary conditions at junctions.}
% source: results/layer0.csv, results/layer0_gci.csv (cone and cone_sym_axis wall_cp true_error_pct: 0.188 and 0.748, a factor of about four)
The tip-junction sensitivity of the cone (a factor of four in the finest-grid wall pressure error) was found in a two-dimensional
structured setup.
It shows that a junction of two boundary types can dominate an error that a grid study would attribute to the
resolution; it does not show that the product's three-dimensional meshes are affected.

\paragraph{Grids the product can build.}
% source: results/gci_existing.csv (header and basic_finner_mesh_check rows)
Because the product fixes the largest cell size, refinement through its settings changes only the near-wall region.
Every GCI of such a family depends on an arbitrary choice of the grid-size definition, and on the Basic Finner the
orders of the two triplets disagree under both.
A reliable GCI of the product would need a mesh control that scales the whole field together with the wall size;
that is a change to the product and outside the scope of this study.

\paragraph{Reference data.}
% source: results/cone10_reference_check.csv; Celik et al. Table 1 cell: tests/test_gci.py (not a results CSV)
Two published verification references each contain an internally inconsistent entry: the NPARC 10$^\circ$ cone table
(Section~\ref{sec:cone10}) and one cell of Table~1 of \citet{celik2008}.
Reporting them is not a criticism of either; it matters for what can be used as a target, and it is why the
reference values here are computed and checked against independent tables.

\paragraph{Limits of this study.}
Layer~0 verifies SU2 with the product's settings and not the product's pipeline.
The laminar plate verifies a configuration the product does not write.
The turbulent-plate data come from a configuration close to but not identical to the product's.
The 2-D boundary-junction and entropy findings are untested on the three-dimensional meshes.
The comparison of Spalart--Allmaras with a second turbulence model was not performed, and the product cannot run one.

%=====================================================================
\section{Conclusions}
\label{sec:conclusions}

The conclusions supported at this stage are the following.
% source: results/layer0.csv, results/layer0_gci.csv, results/cone10_reference_check.csv, results/cone_shock_angle_sources.csv, results/gci_existing.csv, results/temporal.csv (every number in the list); figures/validation_bands.svg (item 6)
\begin{enumerate}
\item With the product's numerical settings, SU2 reproduces exact inviscid solutions to within 0.2\,\% on the surface
  pressure and the shock angle at 257$\times$193 nodes in the primary setups (0.75\,\% on the cone with a symmetry-axis tip
  junction), while wall quantities derived from the Mach number retain errors of 0.8 to 3.6\,\% that decrease slowly or not
  at all.
\item A GCI bounded the true finest-grid error for 5 of 7 inviscid results and for 2 of 5 laminar-plate results, the latter only at
  sign crossings (and the five plate rows are not independent); neither the observed order nor the stability of the order across windows
  detects both failures.
\item The NPARC 10$^\circ$ cone table is internally inconsistent and disagrees with independent tables in every entry except
  the surface Mach number; those tables support a shock angle of 26.74$^\circ$ and not 27.18$^\circ$.
\item Grids built through the product's settings cannot be refined similarly; the four-level Basic Finner Euler study
  does not establish a reliable GCI, and the published RANS sample-rocket grid study fails the convergence gate for every
  condition for which it was tested.
\item The flight integrator shows no convergence order on the sample rocket, attributed on the evidence of controls
  mainly to the rail model, with a second contribution not excluded.
\item Free-flight comparison of the Euler Basic Finner cannot attribute the difference from measurement to model form,
  because no established numerical band exists.
\end{enumerate}
% PENDING Task 18: conclusion on the pipeline error (layer 0 against layer 1)
% PENDING Task 14: conclusion on the sample-rocket Euler grid convergence and d'Alembert check
% PENDING Task 15: conclusion on the Basic Finner RANS grid convergence and the validation overlap
% PENDING Task 16: conclusion on iterative error against the grid uncertainty
% PENDING Task 20: conclusion on Euler against RANS model uncertainty

\section*{Data availability}
All numbers in this paper come from the result files under \texttt{research/nasa\_vv/results/}
(and the committed figure under \texttt{research/nasa\_vv/figures/}) of the IgnisYeet repository,
except where a comment in the source names another committed file.

\bibliographystyle{plainnat}
\bibliography{references}

\end{document}
